\songtitle{Sa fiamma}

\tag*{2024}\tag{italian-folk-pop}\tag{thomas-entourage}\tag{sylvia-lai}\tag{campidanese-lyrics}

\begin{notes}
Based on the \emph{Thomas’ Entourage} fictional universe by Carlos Thompson.\\
Character: Sylvia Lai.\\
Co-written by Carlos Thompson and Claude. %
High interaction.
\end{notes}

\begin{style}
Italian pop, Sardinian folk influences, slow dream-pop. %
A steel-string acoustic guitar plays a steady eighth-note strumming pattern with a bright, percussive attack. %
A clean electric bass guitar follows the root notes of the chord progression with a warm, rounded tone. %
A female vocalist sings in a conversational, mezzo-soprano with a reflective, warm, bittersweet, resolved production style. %
The tempo is 72 BPM in 4/4 time, in the key of G major. %
The arrangement is sparse, focusing on the interplay between the rhythmic acoustic guitar and the vocal melody.
\end{style}

\begin{lyrics}

\songpart{Intro}
\annotation{bright acoustic guitar strumming, warm electric bass}

\songpart{Verse 1}
\annotation{female vocals}\\
Greca e Alessandro ad Elmas con la faccia —\\
nove mesi trattenuti in pubblico.\\
Le domande erano tutte sui capelli.
\annotation{Spoken word}\\
Giusto. %
I capelli sono quello che riescono a vedere.

\songpart{Verse 2}
Roberto sottovoce, una domanda sola —\\
gliene ho dette quattro in campidanese.\\
Tre ore di silenzio, poi è venuto a scusarsi.
\annotation{Spoken word}\\
Alcune cose funzionano meglio in campidanese.

\songpart{Pre-Chorus — Tallinn}
\annotation{Whisper}\\
Sono partita a marzo per due settimane —\\
il mondo ha chiuso i cancelli mentre ero dentro.\\
Ho scelto il peccato con gli occhi aperti —\\
e adesso lo porto senza chiedere perdono.

\songpart{Chorus}
\annotation{Joyful}\\
Sono tornata dove tutto è cominciato,\\
la villa, il fuoco, le voci che conosco.\\
Non sono piú quella che è partita a marzo —\\
ma il Truncu brucia uguale, e deve essere così.
\annotation{Spoken word}\\
Sa fiamma est sa fiamma (est sa fiamma)\\
cambia chi la guarda, non lei. %
(non lei)

\songpart{Verse 3}
Tutti i cugini, i mariti, i bambini che corrono,\\
la nonna che conta le teste intorno al fuoco.\\
La fiamma prende, il legno scoppietta, siamo tutti lì.
\annotation{Spoken word}\\
È ogni anno uguale e ogni anno necessario.

\songpart{Pre-Chorus — Bogotá}
\annotation{Whisper}\\
C'era un appartamento che ho chiamato mio\\
prima ancora di sapere se restavo.\\
Ho deciso il colore senza sapere dove farlo —\\
quando decido una cosa è già decisa.

\songpart{Chorus}
\annotation{Joyful}\\
Sono tornata dove tutto è cominciato,\\
la villa, il fuoco, le voci che conosco.\\
Non sono piú quella che è partita a marzo —\\
ma il Truncu brucia uguale, e deve essere così.
\annotation{Spoken word}\\
Sa fiamma est sa fiamma (est sa fiamma)\\
cambia chi la guarda, non lei. %
(non lei)

\songpart{Bridge}
\annotation{Vulnerable}\\
Quello che è successo in mezzo non si riassume:\\
una città di vetro, uno studio a Queens,\\
qualcuno che mi ha guardata senza costruire una storia.\\
Thomas era la strada — io ero la destinazione.\\
Avevo sedici anni, qualcos'altro adesso —
\annotation{Spoken word}\\
non so ancora il nome, ma lo riconosco.

\songpart{Verse 4}
Le foto di Gabi aperte sullo schermo —\\
un occhio quasi dorato, l'altro sotto i capelli teal.\\
Si contano le lentiggini, i puntini sul naso.
\annotation{Spoken word}\\
A volte li odio. %
In quella foto sono esattamente dove devono essere.

\songpart{Pre-Chorus — New York}
\annotation{Whisper}\\
Non quella che aspettava qualcuno che arrivasse,\\
non quella che chiedeva dove apparteneva.\\
Questa ragazza nello schermo la conosco —\\
potrei innamorarmi di me stessa.

\songpart{Chorus}
\annotation{Joyful}\\
Sono tornata dove tutto è cominciato,\\
la villa, il fuoco, le voci che conosco.\\
Non sono piú quella che è partita a marzo —\\
ma il Truncu brucia uguale, e deve essere così.
\annotation{Spoken word}\\
Sa fiamma est sa fiamma (est sa fiamma)\\
cambia chi la guarda, non lei. %
(non lei)

\annotation{Final Chorus — Campidanese, sparse arrangement, voice forward}\\
Seu torraa in dói chi naschìu,\\
sa villa, su fogu, is boghes chi connòsciu.\\
No seu prus issa chi partìu in màrtu —\\
ma su Truncu brusat sèmpri, e depit èssiri accussì.
\annotation{Spoken word}\\
Sa fiamma est sa fiamma —\\
mùdat chini la tzerriat, no issa.

\songpart{Outro — instrumental fade, acoustic guitar, single piano line}

\end{lyrics}

